% Self-contained report. Compile with pdfLaTeX in Overleaf (twice for contents).
\documentclass[12pt,a4paper]{article}
\usepackage[T1]{fontenc}
\usepackage[utf8]{inputenc}
\usepackage[english]{babel}
\usepackage[margin=2.5cm]{geometry}
\usepackage{xcolor}
\usepackage{listings}
\usepackage{booktabs}
\usepackage{microtype}
\usepackage[hidelinks]{hyperref}
\definecolor{listinggray}{RGB}{247,247,247}
\lstset{basicstyle=\ttfamily\footnotesize,backgroundcolor=\color{listinggray},
 frame=single,rulecolor=\color{gray},breaklines=true,breakatwhitespace=false,
 columns=fullflexible,keepspaces=true,showstringspaces=false,tabsize=4,
 captionpos=b,aboveskip=10pt,belowskip=10pt}
\setlength{\emergencystretch}{3em}
\begin{document}
\begin{titlepage}
\centering
\textbf{Ministry of Education, Culture and Research of the Republic of Moldova}\par\medskip
\textbf{Technical University of Moldova}\par\medskip
\textbf{Faculty of Computers, Informatics and Microelectronics}\par\medskip
\textbf{Department of Software Engineering and Automatic}\par
\vspace{2.5cm}
\textbf{ANDREI ELCISCEV}\par
\textbf{Group FAF-243}\par
\vspace{2cm}
{\Large\textbf{LABORATORY WORK No. 11}\par}
\medskip
{\Large\textbf{THE RACE AND THE SEMAPHORE}\par}
\vspace{2cm}
\begin{flushright}
\textbf{Checked by:}\\
Rei\c{t}man Patricia
\end{flushright}
\vfill
\textbf{Chi\c{s}in\u{a}u, 2026}
\end{titlepage}
\pagenumbering{arabic}
\tableofcontents
\newpage
\section{Introduction}
Two POSIX threads each increment a shared counter one million times. The experiments
compare an unprotected counter with a mutex and a semaphore, then use semaphores to
signal that a message is ready and to limit occupancy in a simulated parking lot.

The programs ran in the Linux VM under \texttt{/home/user/os-lab11}. All five programs
were compiled with \texttt{gcc -Wall -pthread} at the compiler's default optimization
level, with no \texttt{-O2}. The listings show actual commands and output. The altered
initial values were tested in temporary copies of the source files.

\begin{lstlisting}[caption={Linux VM and compiler}]
$ pwd
/home/user/os-lab11
$ uname -sr
Linux 6.14.0-37-generic
$ nproc
3
$ gcc -dumpfullversion
13.3.0
\end{lstlisting}

\section{Part 1: Observe the Race}

\begin{lstlisting}[caption={Three trials without synchronization}]
$ gcc -Wall -pthread -o race race.c
$ ./race
counter = 1675785 (expected 2000000)
$ ./race
counter = 1925845 (expected 2000000)
$ ./race
counter = 1205596 (expected 2000000)
\end{lstlisting}

\paragraph{Observe 1.} Incrementing the counter consists of loading its value,
adding one, and storing the result. These steps do not form an atomic operation.

\paragraph{Observe 2.} Starting from 10, both threads can load 10, compute 11,
and store 11. Two increments have run, but one store overwrote the other update,
leaving 11 instead of 12.

\paragraph{Observe 3.} Scheduling and execution overlap vary across trials, so
the workers can lose a different number of updates. This is a C data race and has
undefined behavior. All three measured trials fell short of 2,000,000.

\section{Part 2: Protect the Counter}
\subsection{Mutex and One-Permit Semaphore}
A mutex encloses the read, addition, and store in one critical section. The
semaphore version also admits one thread at a time when initialized to one.
\texttt{sem\_wait} acquires the permit and \texttt{sem\_post} returns it.

\begin{lstlisting}[caption={Three correctness trials per lock}]
$ gcc -Wall -pthread -o mutex mutex.c
$ ./mutex
counter = 2000000 (expected 2000000)
$ ./mutex
counter = 2000000 (expected 2000000)
$ ./mutex
counter = 2000000 (expected 2000000)
$ gcc -Wall -pthread -o sem sem.c
$ ./sem
counter = 2000000 (expected 2000000)
$ ./sem
counter = 2000000 (expected 2000000)
$ ./sem
counter = 2000000 (expected 2000000)
\end{lstlisting}

\paragraph{Observe 1.} Both protected versions reported 2,000,000 in all three
requested trials and in the three additional timed trials.

\subsection{Semaphore Initially Set to Two}

\begin{lstlisting}[caption={Two permits admit both workers}]
$ sed 's/sem_init(&s, 0, 1)/sem_init(\&s, 0, 2)/' sem.c > sem2.c
$ grep -n sem_init sem2.c
23:    sem_init(&s, 0, 2);         /* 0 = shared between threads, 1 = start value */
$ gcc -Wall -pthread -o sem2 sem2.c
$ ./sem2
counter = 1999341 (expected 2000000)
$ ./sem2
counter = 1947310 (expected 2000000)
$ ./sem2
counter = 1994467 (expected 2000000)
\end{lstlisting}

\paragraph{Observe 2.} Both workers can acquire a permit and enter the critical
section together. The read--modify--write race returns, and all three trials lose
increments. The submitted \texttt{sem.c} uses one permit.

\subsection{The Cost of Synchronization}

\begin{lstlisting}[caption={Three wall-time measurements per version}]
$ time ./race
counter = 1832162 (expected 2000000)
real 0.015 s | user 0.015 s | sys 0.002 s
$ time ./mutex
counter = 2000000 (expected 2000000)
real 0.046 s | user 0.060 s | sys 0.017 s
$ time ./sem
counter = 2000000 (expected 2000000)
real 0.074 s | user 0.087 s | sys 0.044 s
$ time ./race
counter = 1269920 (expected 2000000)
real 0.010 s | user 0.014 s | sys 0.002 s
$ time ./mutex
counter = 2000000 (expected 2000000)
real 0.058 s | user 0.055 s | sys 0.016 s
$ time ./sem
counter = 2000000 (expected 2000000)
real 0.072 s | user 0.072 s | sys 0.018 s
$ time ./race
counter = 1269203 (expected 2000000)
real 0.017 s | user 0.016 s | sys 0.006 s
$ time ./mutex
counter = 2000000 (expected 2000000)
real 0.069 s | user 0.068 s | sys 0.015 s
$ time ./sem
counter = 2000000 (expected 2000000)
real 0.072 s | user 0.081 s | sys 0.050 s
\end{lstlisting}

\paragraph{Observe 3.} Acquiring and releasing a lock on every increment adds
synchronization operations. When workers contend, the shared lock state also
causes cache-coherence traffic and can require waits and wakeups.
The median wall times in this VM were 0.015 s for the race, 0.058 s for the mutex, and 0.072 s for the semaphore.

\section{Part 3: Use a Semaphore as a Signal}
The reader waits on \texttt{ready}. With an initial value of zero, the writer
stores the message and posts before the reader can acquire a permit and read it.

\begin{lstlisting}[caption={Signal experiment with initial values zero and one}]
$ gcc -Wall -pthread -o order order.c
$ ./order
B: waiting for the message...
A: message written
B: I read "hello"
$ sed 's/sem_init(&ready, 0, 0)/sem_init(\&ready, 0, 1)/' order.c > order1.c
$ grep -n sem_init order1.c
30:    sem_init(&ready, 0, 1);
$ gcc -Wall -pthread -o order1 order1.c
$ ./order1
B: waiting for the message...
B: I read ""
A: message written
\end{lstlisting}

\paragraph{Observe 1.} With an initial value of one, the reader consumed the
initial permit and printed an empty string before the writer ran its store. The
global array starts with zero bytes. This variant removes the synchronization
ordering between the writer's store and the reader's load, so the empty string
describes this run.

\paragraph{Observe 2.} A semaphore initialized to zero makes a waiter wait for a
post to supply a permit. The submitted version starts at zero and the reader
printed \texttt{hello} after the writer's message.

\section{Part 4: Count Free Parking Spaces}
Each of the six car threads acquires a permit before entering and returns it
after leaving. The mutex protects the occupancy count and the associated print.

\begin{lstlisting}[caption={Parking with three places}]
$ gcc -Wall -pthread -o parking parking.c
$ ./parking
car 1 enters  (inside: 1)
car 3 enters  (inside: 2)
car 2 enters  (inside: 3)
car 2 leaves  (inside: 2)
car 3 leaves  (inside: 1)
car 5 enters  (inside: 2)
car 4 enters  (inside: 3)
car 1 leaves  (inside: 2)
car 6 enters  (inside: 3)
car 4 leaves  (inside: 2)
car 6 leaves  (inside: 1)
car 5 leaves  (inside: 0)
\end{lstlisting}

\begin{lstlisting}[caption={Parking with one place}]
$ sed 's/^#define PLACES 3$/#define PLACES 1/' parking.c > parking1.c
$ grep -n #define PLACES parking1.c
9:#define PLACES 1
$ gcc -Wall -pthread -o parking1 parking1.c
$ ./parking1
car 1 enters  (inside: 1)
car 1 leaves  (inside: 0)
car 2 enters  (inside: 1)
car 2 leaves  (inside: 0)
car 4 enters  (inside: 1)
car 4 leaves  (inside: 0)
car 5 enters  (inside: 1)
car 5 leaves  (inside: 0)
car 6 enters  (inside: 1)
car 6 leaves  (inside: 0)
car 3 enters  (inside: 1)
car 3 leaves  (inside: 0)
\end{lstlisting}

\paragraph{Observe 1.} The maximum recorded occupancy was three.
\texttt{sem\_init(\&places, 0, PLACES)} provides three permits, and
\texttt{sem\_wait(\&places)} gates entry. A car releases its permit with
\texttt{sem\_post(\&places)} after leaving.

\paragraph{Observe 2.} With one place, the semaphore functions as a binary
semaphore for the slot. The trace shows one car at a time. All six cars entered
and left once in each experiment, and both runs ended with zero cars inside.
The submitted version uses three places.

\section{Conclusion}
The mutex and the semaphore with one permit prevented lost increments. Starting
the message semaphore at zero enforced the ordering between writer and reader.
Starting the parking semaphore at three bounded concurrent occupancy. Each use
depends on when a thread acquires or returns a permit and how many permits are
available initially.

\appendix
\section{Submitted Source Code}

\begin{lstlisting}[caption={race.c},language=C]
/* race.c - two threads add to the same counter, with no protection.
 * Build:  gcc -Wall -pthread -o race race.c      (no -O2: keep counter++ in memory)
 * Run:    ./race   (several times)
 */
#include <stdio.h>
#include <pthread.h>

#define N 1000000

long counter = 0;                 /* shared by both threads */

void *add(void *arg) {
    for (int i = 0; i < N; i++)
        counter++;                /* read, add 1, write back: three steps */
    return NULL;
}

int main(void) {
    pthread_t a, b;
    pthread_create(&a, NULL, add, NULL);
    pthread_create(&b, NULL, add, NULL);
    pthread_join(a, NULL);
    pthread_join(b, NULL);
    printf("counter = %ld (expected %d)\n", counter, 2 * N);
    return 0;
}
\end{lstlisting}

\begin{lstlisting}[caption={mutex.c},language=C]
/* mutex.c - the same counter, now protected by a mutex.
 * Build:  gcc -Wall -pthread -o mutex mutex.c
 */
#include <stdio.h>
#include <pthread.h>

#define N 1000000

long counter = 0;
pthread_mutex_t lock = PTHREAD_MUTEX_INITIALIZER;

void *add(void *arg) {
    for (int i = 0; i < N; i++) {
        pthread_mutex_lock(&lock);      /* only one thread past this line */
        counter++;                      /* the critical section */
        pthread_mutex_unlock(&lock);    /* let the other one in */
    }
    return NULL;
}

int main(void) {
    pthread_t a, b;
    pthread_create(&a, NULL, add, NULL);
    pthread_create(&b, NULL, add, NULL);
    pthread_join(a, NULL);
    pthread_join(b, NULL);
    printf("counter = %ld (expected %d)\n", counter, 2 * N);
    return 0;
}
\end{lstlisting}

\begin{lstlisting}[caption={sem.c},language=C]
/* sem.c - the same counter, protected by a semaphore that starts at 1.
 * Build:  gcc -Wall -pthread -o sem sem.c
 */
#include <stdio.h>
#include <pthread.h>
#include <semaphore.h>

#define N 1000000

long counter = 0;
sem_t s;

void *add(void *arg) {
    for (int i = 0; i < N; i++) {
        sem_wait(&s);           /* 1 -> 0: I go in; at 0: wait */
        counter++;
        sem_post(&s);           /* 0 -> 1: the next one may go in */
    }
    return NULL;
}

int main(void) {
    sem_init(&s, 0, 1);         /* 0 = shared between threads, 1 = start value */
    pthread_t a, b;
    pthread_create(&a, NULL, add, NULL);
    pthread_create(&b, NULL, add, NULL);
    pthread_join(a, NULL);
    pthread_join(b, NULL);
    printf("counter = %ld (expected %d)\n", counter, 2 * N);
    sem_destroy(&s);
    return 0;
}
\end{lstlisting}

\begin{lstlisting}[caption={order.c},language=C]
/* order.c - thread B must wait until thread A has written the message.
 * A semaphore that starts at 0 is a signal: sem_wait blocks until someone posts.
 * Build:  gcc -Wall -pthread -o order order.c
 */
#include <stdio.h>
#include <string.h>
#include <unistd.h>
#include <pthread.h>
#include <semaphore.h>

char message[32];
sem_t ready;                    /* starts at 0: nothing is ready yet */

void *writer(void *arg) {
    sleep(1);                   /* pretend the work takes a while */
    strcpy(message, "hello");
    printf("A: message written\n");
    sem_post(&ready);           /* 0 -> 1: wake up B */
    return NULL;
}

void *reader(void *arg) {
    printf("B: waiting for the message...\n");
    sem_wait(&ready);           /* blocks while the value is 0 */
    printf("B: I read \"%s\"\n", message);
    return NULL;
}

int main(void) {
    sem_init(&ready, 0, 0);
    pthread_t a, b;
    pthread_create(&b, NULL, reader, NULL);   /* B starts first, on purpose */
    pthread_create(&a, NULL, writer, NULL);
    pthread_join(a, NULL);
    pthread_join(b, NULL);
    sem_destroy(&ready);
    return 0;
}
\end{lstlisting}

\begin{lstlisting}[caption={parking.c},language=C]
/* parking.c - 6 cars, 3 places. A semaphore that starts at 3 counts the free places.
 * Build:  gcc -Wall -pthread -o parking parking.c
 */
#include <stdio.h>
#include <unistd.h>
#include <pthread.h>
#include <semaphore.h>

#define PLACES 3
#define CARS   6

sem_t places;                                   /* free places */
pthread_mutex_t lock = PTHREAD_MUTEX_INITIALIZER;
int inside = 0;                                 /* only for printing */

void *car(void *arg) {
    long id = (long)arg;
    sem_wait(&places);                          /* take a place, or wait */
    pthread_mutex_lock(&lock);
    inside++;
    printf("car %ld enters  (inside: %d)\n", id, inside);
    pthread_mutex_unlock(&lock);

    sleep(1);                                   /* parked */

    pthread_mutex_lock(&lock);
    inside--;
    printf("car %ld leaves  (inside: %d)\n", id, inside);
    pthread_mutex_unlock(&lock);
    sem_post(&places);                          /* give the place back */
    return NULL;
}

int main(void) {
    sem_init(&places, 0, PLACES);
    pthread_t t[CARS];
    for (long i = 0; i < CARS; i++)
        pthread_create(&t[i], NULL, car, (void *)(i + 1));
    for (int i = 0; i < CARS; i++)
        pthread_join(t[i], NULL);
    sem_destroy(&places);
    return 0;
}
\end{lstlisting}

\end{document}
